\chapter{Computational Appendix: Reproducibility and Data Transparency}

\section{Principal Component Analysis (PCA): Objective Distance Clustering}
This appendix provides the Python implementation of the unsupervised clustering utilized in Chapter 4. Unlike previous models, this approach uses normalized Levenshtein distance (edit distance) on raw primary data strings to ensure a non-circular proof of structural signals.

\begin{verbatim}
import numpy as np
from sklearn.decomposition import PCA

def levenshtein_dist(s1, s2):
    """Calculates the normalized Levenshtein distance."""
    if len(s1) < len(s2): return levenshtein_dist(s2, s1)
    if len(s2) == 0: return len(s1)
    previous_row = range(len(s2) + 1)
    for i, c1 in enumerate(s1):
        current_row = [i + 1]
        for j, c2 in enumerate(s2):
            insertions = previous_row[j + 1] + 1
            deletions = current_row[j] + 1
            substitutions = previous_row[j] + (c1 != c2)
            current_row.append(min(insertions, deletions, substitutions))
        previous_row = current_row
    return previous_row[-1] / max(len(s1), len(s2))

def run_unsupervised_discovery(all_aligned_sets):
    N = len(all_aligned_sets)
    manifold = np.zeros((N, N))
    for i in range(N):
        for j in range(N):
            d = 0
            for k in range(3): # Vedic, Tamil, Tibetan
                d += levenshtein_dist(all_aligned_sets[i][k], 
                                      all_aligned_sets[j][k])
            manifold[i, j] = d / 3
            
    pca = PCA(n_components=2)
    latent_space = pca.fit_transform(manifold)
    return pca.explained_variance_ratio_[0]
\end{verbatim}

\section{Hidden Markov Model (HMM): Structural Null Model Baseline}
The following script documents the HMM architecture used for the Zero-Shot predictive validation in Chapter 5. The model is tasked with predicting the transitions of a held-out Harappan dataset against a standard SOV baseline.

\begin{verbatim}
from hmmlearn import hmm
import numpy as np

# Model configuration
n_states = 4
model = hmm.CategoricalHMM(n_components=n_states, n_iter=200, random_state=42)

# Structural SOV Null Model (Transition Matrix Baseline)
# Represents a standard right-branching SOV syntax
A_sov = np.array([
    [0.1, 0.7, 0.1, 0.1], # Agent -> Target (high)
    [0.1, 0.1, 0.7, 0.1], # Target -> Verb (high)
    [0.1, 0.1, 0.1, 0.7], # Verb -> Terminal (high)
    [0.7, 0.1, 0.1, 0.1]  # Terminal -> New Agent (high)
])

# Training protocol:
# 1. Purge all semantic and lexical labels from corpus.
# 2. Split dataset 80/20 into training and held-out validation sets.
# 3. Apply Laplacian additive smoothing (epsilon = 1e-6).
# 4. Evaluate predictive Log-Likelihood against the SOV Null Matrix.
\end{verbatim}

\section{Indus Valley Script: Cleaned Dataset Sample (n=50)}
A representative sample of the cleaned structural data used for the HMM transitions. Sequences represent numeric sign identifiers from the Mahadevan concordance, purged of all null delimiters.

\begin{table}[h]
\centering
\begin{tabular}{ll|ll} \toprule
\textbf{Seal ID} & \textbf{Sequence (M-codes)} & \textbf{Seal ID} & \textbf{Sequence (M-codes)} \\ \midrule
1.1 & 410, 017 & 106.1 & 740, 235 \\
2.1 & 410, 017 & 110.1 & 740, 017 \\
3.1 & 405, 017 & 115.1 & 390, 817 \\
5.1 & 740, 235 & 120.1 & 740, 061 \\
6.1 & 740, 390, 590 & 125.1 & 390, 004 \\
7.1 & 368, 390, 125, 033 & 130.1 & 002, 817 \\
8.1 & 740, 904, 033, 705, 235 & 135.1 & 236, 002, 308, 072 \\
9.1 & 740, 017, 495, 235, 741, 066 & 140.1 & 740, 803 \\
10.1 & 156, 460, 510 & 145.1 & 426, 003, 060, 692 \\
11.1 & 072, 170, 740, 840, 013 & 150.1 & 374, 308, 350 \\
13.1 & 740, 440, 503, 002, 861 & 155.1 & 740, 017 \\
16.1 & 390, 415 & 160.1 & 700, 017 \\
17.1 & 390, 034 & 165.1 & 017, 390 \\
18.1 & 405, 844, 032, 840, 002, 861 & 170.1 & 740, 235, 861 \\
19.1 & 226, 003, 002, 297, 350, 125, 413 & 175.1 & 368, 017 \\
20.1 & 390, 717, 061, 002, 368, 634 & 180.1 & 740, 900, 017 \\
22.1 & 527, 554, 298, 368, 634 & 185.1 & 390, 004, 817 \\
23.1 & 740, 798 & 190.1 & 740, 061 \\
26.1 & 520, 001 & 195.1 & 002, 817 \\
27.1 & 090, 740, 220, 002, 368, 550, 821 & 200.1 & 390, 415 \\
28.1 & 220, 014 & 205.1 & 740, 235 \\
29.1 & 740, 797 & 210.1 & 017, 368 \\
30.1 & 002, 817 & 215.1 & 740, 590 \\
31.1 & 002, 817 & 220.1 & 390, 705 \\
32.1 & 740, 061 & 225.1 & 740, 235 \\
33.1 & 390, 004 & 230.1 & 368, 017 \\ \bottomrule
\end{tabular}
\caption{Representative 50-Seal Sample from the Cleaned Harappan Corpus.}
\end{table}

